% !TeX program = xelatex
% !BIB program = biber
\documentclass[light]{lutbeamer}
\usepackage{pgfpages}
\usepackage{ragged2e}
\usepackage{array}
\usepackage{booktabs}
\usepackage{amssymb}
\usepackage{fontawesome5}
\usepackage{tikz}
\usepackage{pgfplots}
\usetikzlibrary{arrows.meta,positioning,calc}
\pgfplotsset{compat=1.18}
\addbibresource{references.bib}

\setdepartment{OSTIM Teknik Üniversitesi}
\institute{}
\author[Mühendislik Fakültesi]{\\
Prof.~Dr.~Yalçın ATA \\
Prof.~Dr.~Arif DEMİR \\
Dr.~Öğr.~Üyesi~Muhammed ELMNEFİ \\
Dr.~Öğr.~Üyesi~Haydar KILIÇ \\
Dr.~Öğr.~Üyesi~Emel GÜVEN \\
Dr.~Öğr.~Üyesi~Ufuk ASIL \\
Öğr.~Gör.~Sema ÇİFTÇİ
}
\title{MATH 204 -- Olasılık ve İstatistik}
\subtitle{Hafta 14: Vize Öncesi Zor Düzey 20 Soru}
\date{2025--2026 Bahar Dönemi}

\definecolor{probExampleGreen}{HTML}{2E7D32}
\definecolor{probInfoGray}{HTML}{4B5563}
\definecolor{probResultOrange}{HTML}{FF8621}

\setbeamercolor{block title example}{fg=white,bg=probExampleGreen}
\setbeamercolor{block body example}{fg=black,bg=probExampleGreen!8}

\newenvironment{infoblock}[1]
{
    \setbeamercolor{block title}{fg=white,bg=probInfoGray}
    \setbeamercolor{block body}{fg=black,bg=probInfoGray!10}
    \begin{block}{#1}
}
{
    \end{block}
}

\newenvironment{noteblock}[1]
{
    \begin{infoblock}{#1}
}
{
    \end{infoblock}
}

\newenvironment{resultblock}[1]
{
    \setbeamercolor{block title}{fg=white,bg=probResultOrange}
    \setbeamercolor{block body}{fg=black,bg=probResultOrange!12}
    \begin{block}{#1}
}
{
    \end{block}
}

\newcommand{\solutionghost}[1]{{\color{red!45!black}#1}}
\newcommand{\solutiontext}[1]{{\color{red!85!black}#1}}
\newcommand{\ghoststep}[2]{\textbf{\solutionghost{Adım #1.}}~\solutionghost{#2}\par\vspace{0.08em}}
\newcommand{\ghostmath}[1]{%
\begingroup
\setlength{\abovedisplayskip}{4pt}
\setlength{\belowdisplayskip}{4pt}
\setlength{\abovedisplayshortskip}{3pt}
\setlength{\belowdisplayshortskip}{3pt}
\[
\color{red!45!black}
#1
\]
\endgroup
\vspace{0.05em}}
\newcommand{\solutionstep}[2]{\textbf{\solutiontext{Adım #1.}}~\solutiontext{#2}\par\vspace{0.08em}}
\newcommand{\solutionmath}[1]{%
\begingroup
\setlength{\abovedisplayskip}{4pt}
\setlength{\belowdisplayskip}{4pt}
\setlength{\abovedisplayshortskip}{3pt}
\setlength{\belowdisplayshortskip}{3pt}
\[
\color{red!85!black}
#1
\]
\endgroup
\vspace{0.05em}}
\newcommand{\solutionfinal}[1]{%
\vfill
{\color{probResultOrange}\rule{\linewidth}{0.8pt}\par}
\vspace{0.2em}
\noindent\textbf{\solutiontext{Sonuç:}}~\solutiontext{#1}}
\newcommand{\solutionsetup}[3]{%
\noindent\solutiontext{\textbf{Soru türü:} #1 \quad \textbf{Araç:} #2 \quad \textbf{Neden:} #3}\par
\vspace{0.2em}}
\newcommand{\solutionmistake}[1]{%
\noindent{\color{red!75!black}\textbf{Sık hata:} #1}\par
\vspace{0.15em}}

\setcounter{tocdepth}{2}

\AtBeginSection[]{
  \begin{frame}{Sunum Planı}
    \scriptsize
    \begin{columns}[T]
      \begin{column}{0.48\textwidth}
        \tableofcontents[sections={1-3}, currentsection]
      \end{column}
      \begin{column}{0.48\textwidth}
        \tableofcontents[sections={4-7}, currentsection]
      \end{column}
    \end{columns}
  \end{frame}
}

\begin{document}

\setbeamertemplate{navigation symbols}{}
\setbeamertemplate{footline}{
  \begin{beamercolorbox}[wd=\paperwidth, ht=10pt, dp=1pt]{footline}
    \usebeamerfont{footline}
    \begin{tikzpicture}[remember picture,overlay]
      \node[anchor=south west, xshift=10mm, yshift=.4mm]
      at (current page.south west)
      {\insertshortauthor,~\department};
    \end{tikzpicture}
    \begin{tikzpicture}[remember picture,overlay]
      \node[anchor=south east, xshift=-5mm, yshift=0.5mm]
      at (current page.south east)
      {\insertframenumber \,/\, \inserttotalframenumber};
    \end{tikzpicture}
  \end{beamercolorbox}
}

\begin{frame}<article:0>[plain,noframenumbering]
  \begin{tikzpicture}[remember picture,overlay]
    \node[anchor=center, at=(current page.center), inner sep=0pt] {
      \includegraphics[width=\paperwidth, height=\paperheight]{logos/background_figure.jpg}
    };
  \end{tikzpicture}
\end{frame}

{
  \setbeamertemplate{footline}{
    \begin{beamercolorbox}[wd=\paperwidth, ht=10pt, dp=1pt]{footline}
      \usebeamerfont{footline}
      \begin{tikzpicture}[remember picture,overlay]
        \node[anchor=south west, xshift=10mm, yshift=.4mm, align=left]
        at (current page.south west)
        {\tiny Şablon bağlantısı: \href{https://github.com/ufukasia/Ostim-Teknik--niversitesi-Yaz-l-m-M-hendisli-i---Lutbeamer-egitim-i-eri-i-haz-rlama-}{GitHub}.};
      \end{tikzpicture}
      \begin{tikzpicture}[remember picture,overlay]
        \node[anchor=south east, xshift=-5mm, yshift=0.5mm]
        at (current page.south east)
        {\insertframenumber \,/\, \inserttotalframenumber};
      \end{tikzpicture}
    \end{beamercolorbox}
  }
  \setbeamerfont{title}{size=\Large}
  \setbeamerfont{author}{size=\scriptsize}
  \setbeamerfont{subtitle}{size=\small}
  \setbeamerfont{date}{size=\footnotesize}
  \maketitle{}
}

\begin{frame}{İçindekiler}
  \scriptsize
  \begin{columns}[T]
    \begin{column}{0.48\textwidth}
      \tableofcontents[sections={1-3}]
    \end{column}
    \begin{column}{0.48\textwidth}
      \tableofcontents[sections={4-7}]
    \end{column}
  \end{columns}
\end{frame}

\section{Hafta Çerçevesi}
\subsection{Hedef}

\begin{frame}[t]{Haftanın Hedefi ve Kapsamı}
\textbf{Amaç}
\begin{itemize}
  \item Yaklaşan vize sınavı öncesinde genel tekrar yapılmasını sağlamak.
  \item Temel olasılık kavramlarını ileri düzey problem çözümleriyle pekiştirmek.
  \item Sınav formatına uygun \textbf{zor düzey} sorular üzerinden analiz yeteneğini geliştirmek.
\end{itemize}

\vspace{0.4em}
\textbf{İçerik Dağılımı}
\begin{enumerate}
    \item Örnek Uzay, Olaylar ve Olay Cebiri
    \item Bağımsızlık, Koşullu Olasılık ve Bayes Kuralı
    \item Ayrık ve Sürekli Rassal Değişkenler ile Olasılık Dağılımları
\end{enumerate}
\end{frame}

\subsection{Genel Strateji}
\begin{frame}[t]{Problem Çözüm Stratejileri}
\small
\begin{columns}[T,onlytextwidth]
  \begin{column}{0.48\textwidth}
    \textbf{Olasılık ve Olay Sorularında}
    \begin{itemize}
      \item Önce deneyin işleyişini anlayarak \textbf{örnek uzayı} netleştirin.
      \item Verilenleri ve istenenleri kesişim, birleşim veya tümleyen olarak \textbf{kümelere} dökün.
      \item Koşullu olasılık ve Bayes teoremi için tam olasılık yasasını kurun.
    \end{itemize}
  \end{column}
  \begin{column}{0.48\textwidth}
    \textbf{Dağılım ve Değişken Sorularında}
    \begin{itemize}
      \item Değişkenin ayrık mı yoksa sürekli mi olduğuna baştan karar verin.
      \item Fonksiyonun \textbf{destek (tanım) kümesini} belirleyerek sınırları yazın.
      \item Olasılık fonsiyonu değerleri toplamı/integrali \(1\) olmalıdır ilkesiyle bilinmeyen parametreleri bulun.
    \end{itemize}
  \end{column}
\end{columns}

\vspace{0.6em}
\begin{infoblock}{Problem Akışı}
\textbf{Soru \(\rightarrow\) Soru Durağı \(\rightarrow\) Tam Çözüm} zinciri ile ilerlenerek, çözüm stratejisinin öğrenci tarafından adım adım düşünülmesi hedeflenmektedir.
\end{infoblock}
\end{frame}

\section{Olay Cebiri ve Bağımsızlık}
\subsection{Örnek Uzay ve Kümeler}

\begin{frame}[t]{Soru 1: İki Zar Olayları}
\footnotesize
İki farklı zar aynı anda atılıyor.

\[
A=\{\text{toplam } \ge 9\},\qquad
B=\{\text{en az bir zar }2\},\qquad
C=\{\text{ilk zar }>4\}.
\]

\textbf{İstenen}
\begin{itemize}
  \item[A.] \(A\), \(B\), \(C\) olaylarını elemanlarıyla yazınız.
  \item[B.] \(A\cap B\), \(A\cap C\), \(B\cap C\), \((A\cup B)\cap C\) bulunuz.
  \item[C.] \(A\) ile \(B\) bağımsız mı?
  \item[D.] \(B\) ile \(C\) ayrık mı?
\end{itemize}
\end{frame}

\begin{frame}[t]{Soru 1: Çözüm İçin Durak}
\scriptsize
\solutionghost{
\textbf{Kurulum}
\[
\Omega=\{(i,j): i,j\in\{1,2,3,4,5,6\}\}.
\]
\[
A=\{\text{toplam } \ge 9\},\qquad
B=\{\text{en az bir zar }2\},\qquad
C=\{\text{ilk zar }>4\}.
\]

\textbf{Kendinize sorun}
\begin{itemize}
  \item \(A\)'nın elemanlarını yazarken hangi toplamlar mümkündür?
  \item \(B\cap C\) boş mu, değil mi?
  \item Bağımsızlık için hangi eşitliği test edeceksiniz?
\end{itemize}
}
\end{frame}

\begin{frame}[t,shrink=12]{Soru 1: Çözüm}
\scriptsize
\solutionsetup{olay cebiri ve bağımsızlık}{örnek uzay, küme işlemleri, bağımsızlık testi}{istenenlerin tamamı olayların elemanları ve kesişimleri üzerinden okunur.}
\solutionmistake{Ayrık olmak ile bağımsız olmak aynı şey değildir; boş kesişim çoğu zaman bağımsızlık vermez.}
\solutionstep{1}{Önce olayları açık küme biçiminde yazarak kesişimleri görünür hale getiriyoruz.}
\solutionmath{
\begin{aligned}
A={}&\{(3,6),(4,5),(5,4),(6,3),(4,6),\\
& (5,5),(6,4),(5,6),(6,5),(6,6)\},\\
B={}&\{(2,1),(2,2),(2,3),(2,4),(2,5),(2,6),\\
& (1,2),(3,2),(4,2),(5,2),(6,2)\},\\
C={}&\{(5,j),(6,j):j=1,\ldots,6\}.
\end{aligned}
}
\solutionmath{
\begin{aligned}
A\cap B&=\varnothing,\\
A\cap C&=\{(5,4),(5,5),(5,6),(6,3),(6,4),(6,5),(6,6)\},\\
B\cap C&=\{(5,2),(6,2)\}.
\end{aligned}
}

\solutionstep{2}{Önce \(A\cup B\) kümesini düşünüp sonra bunu \(C\) ile kesiyoruz.}
\solutionmath{
(A\cup B)\cap C=\{(5,2),(5,4),(5,5),(5,6),(6,2),(6,3),(6,4),(6,5),(6,6)\}.
}

\solutionstep{3}{Bağımsızlık şartı olan \(P(A\cap B)=P(A)P(B)\) ve ayrıklık şartı olan \(B\cap C=\varnothing\) eşitliklerini kontrol ediyoruz.}
\solutionmath{
P(A\cap B)=0,\qquad
P(A)P(B)=\frac{10}{36}\cdot\frac{11}{36}>0.
}
\solutionfinal{Bu zar bağlamında \(A\) ile \(B\) birlikte olamadığı için bağımsız değildir; ayrıca \(B\) ile \(C\) aynı atışta birlikte gerçekleşebildiğinden ayrık da değildir.}
\end{frame}

\begin{frame}[t]{Soru 2: Üç Üründe Hasar}
\footnotesize
Bir üretim hattından art arda üç ürün seçiliyor; her ürün hasarlı (\(H\))
veya sağlam (\(S\)) olabilir.

\textbf{İstenen}
\begin{itemize}
  \item[A.] En ayrıntılı örnek uzayı \(\Omega\) ile yazınız.
  \item[B.] \(A\): tam bir hasarlı, \(B\): ilk ürün hasarlı, \(C\): son ürün sağlam.
  \item[C.] Eş olasılıklı temel sonuçlarla \(P(A)\), \(P(B)\), \(P(C)\) bulunuz.
  \item[D.] \(A\) ile \(B\) bağımsız mı?
\end{itemize}
\end{frame}

\begin{frame}[t]{Soru 2: Çözüm İçin Durak}
\small
\solutionghost{
\textbf{Kurulum}
\[
\Omega=\{HHH,HHS,HSH,HSS,SHH,SHS,SSH,SSS\}.
\]

\[
A=\{\text{tam bir hasarlı}\},\quad
B=\{\text{ilk ürün hasarlı}\},\quad
C=\{\text{son ürün sağlam}\}.
\]

\textbf{Kendinize sorun}
\begin{itemize}
  \item Her olay kaç temel sonuç içeriyor?
  \item \(A\cap B\) içinde hangi diziler kalıyor?
  \item Bağımsızlık için hangi iki niceliği karşılaştırmanız gerekiyor?
\end{itemize}
}
\end{frame}

\begin{frame}[t,shrink=10]{Soru 2: Çözüm}
\scriptsize
\solutionsetup{eş olasılıklı örnek uzay ve bağımsızlık}{sayma, eleman sayısı, bağımsızlık testi}{temel sonuçlar eş olasılıklı olduğundan önce olayları saymak gerekir.}
\solutionmistake{Bağımsızlık için yalnızca \(A\cap B\) sayılmaz; mutlaka \(P(A\cap B)=P(A)P(B)\) karşılaştırması yapılır.}
\solutionstep{1}{Önce olayları açıkça yazıp eleman sayılarını buluyoruz; çünkü olasılıklar doğrudan bu sayılardan gelecek.}
\solutionmath{
\Omega=\{HHH,HHS,HSH,HSS,SHH,SHS,SSH,SSS\}.
}
\solutionmath{
A=\{HSS,SHS,SSH\},\quad
B=\{HHH,HHS,HSH,HSS\},\quad
C=\{HHS,HSS,SHS,SSS\}.
}
\solutionmath{
|A|=3,\qquad |B|=4,\qquad |C|=4
\Longrightarrow
P(A)=\frac{3}{8},\quad P(B)=\frac{1}{2},\quad P(C)=\frac{1}{2}.
}

\solutionstep{2}{\(A\) ve \(B\)'nin ortak sonucu yalnızca \(HSS\)'dir.}
\solutionmath{
A\cap B=\{HSS\},\qquad P(A\cap B)=\frac{1}{8}.
}

\solutionstep{3}{İki olayın bağımsız olması için gereken \(P(A\cap B)=P(A)P(B)\) şartını test ediyoruz.}
\solutionmath{
P(A\cap B)=\frac{1}{8},\qquad P(A)P(B)=\frac{3}{8}\cdot\frac{1}{2}=\frac{3}{16}.
}
\solutionfinal{Bu üretim örneğinde “tam bir hasarlı” bilgisi, ilk ürünün hasarlı olmasıyla ilişkili çıktığı için iki olay bağımsız değildir.}
\end{frame}

\begin{frame}[t]{Soru 3: Tolerans ve Görsel Kontrol}
\footnotesize
Bir parçanın tolerans içinde olması olayı \(T\), görsel kontrolden geçmesi olayı \(G\) olsun.
\[
P(T)=0.6,\qquad P(G)=0.5,\qquad P(T\cup G)=0.8.
\]

\textbf{İstenen}
\begin{itemize}
  \item[A.] \(P(T\cap G)\) bulunuz.
  \item[B.] \(P(T\mid G)\) ve \(P(G\mid T)\) bulunuz.
  \item[C.] \(T\) ile \(G\) bağımsız mı?
  \item[D.] \(P(T'\cap G)\) ve \(P(T'\cup G')\) bulunuz.
\end{itemize}
\end{frame}

\begin{frame}[t]{Soru 3: Çözüm İçin Durak}
\small
\solutionghost{
\textbf{Kurulum}
\[
P(T)=0.6,\qquad P(G)=0.5,\qquad P(T\cup G)=0.8.
\]

\textbf{İlk hangi kural?}
\[
P(T\cup G)=P(T)+P(G)-P(T\cap G).
\]

\textbf{Kendinize sorun}
\begin{itemize}
  \item Önce \(P(T\cap G)\) bulunmadan hangi nicelikler hesaplanamaz?
  \item Bağımsızlık için hangi eşitlik test edilecek?
  \item Tamamlayıcı olaylar hangi kısa ilişkiyle yazılabilir?
\end{itemize}
}
\end{frame}

\begin{frame}[t,shrink=6]{Soru 3: Çözüm}
\scriptsize
\solutionsetup{birleşim, koşullu olasılık ve bağımsızlık}{birlik formülü, koşullu olasılık tanımı, tamamlayıcı olaylar}{verilen üç sayıdan ortak olayı çekmeden diğer adımlar kurulamaz.}
\solutionmistake{\(P(A\mid B)\) ile \(P(B\mid A)\) aynı değildir; hangi bilginin paydayı daralttığını açık yazın.}
\solutionstep{1}{Önce birleşim formülüyle ortak olasılığı buluyoruz; çünkü koşullu olasılıklar bu değere dayanır.}
\solutionmath{
P(T\cap G)=P(T)+P(G)-P(T\cup G)=0.6+0.5-0.8=0.3.
}

\solutionstep{2}{Koşullu olasılıkları doğrudan tanımdan yazıyoruz.}
\solutionmath{
P(T\mid G)=\frac{P(T\cap G)}{P(G)}=\frac{0.3}{0.5}=0.6,\qquad
P(G\mid T)=\frac{P(T\cap G)}{P(T)}=\frac{0.3}{0.6}=0.5.
}

\solutionstep{3}{Bağımsızlık için \(P(T\cap G)=P(T)P(G)\) şartını ve tamamlayıcı olayları kontrol ediyoruz.}
\solutionmath{
P(T\cap G)=0.3,\qquad P(T)P(G)=0.6\cdot 0.5=0.3 \implies P(T\cap G)=P(T)P(G).
}
\solutionmath{
P(T'\cap G)=P(G)-P(T\cap G)=0.2,\qquad
P(T'\cup G')=1-P(T\cap G)=0.7.
}
\solutionfinal{Bu kalite kontrol bağlamında tolerans içinde olmak ile görsel kontrolden geçmek bağımsız görünüyor; ancak aynı parça iki koşulu birden sağlayabildiği için olaylar ayrık değildir.}
\end{frame}

\begin{frame}[t]{Soru 4: Üç Bağımsız Proje}
\footnotesize
Birbirinden tamamen bağımsız kurullarda değerlendirilen üç farklı proje için; \(A\) projesinin onaylanması olayı, \(K\) projesinin onaylanması olayı ve \(U\) projesinin onaylanması olayı olsun.
\[
P(A)=P(K)=P(U)=\frac12,\qquad
P(A\cap K)=P(A\cap U)=P(K\cap U)=\frac14,
\]
\[
P(A\cap K\cap U)=\frac18.
\]

\textbf{İstenen}
\begin{itemize}
  \item[A.] İkili bağımsızlığı gösteriniz.
  \item[B.] Tam bağımsızlığı gösteriniz.
  \item[C.] \(P(A\cup K\cup U)\) değerini bulunuz.
\end{itemize}
\end{frame}

\begin{frame}[t]{Soru 4: Çözüm İçin Durak}
\small
\solutionghost{
\textbf{Kurulum}
\[
P(A)=P(K)=P(U)=\frac12,
\]

\[
P(A\cap K)=P(A\cap U)=P(K\cap U)=\frac14,\qquad
P(A\cap K\cap U)=\frac18.
\]

\textbf{İlk hangi araç?}
\[
P(A\cap K)\stackrel{?}{=}P(A)P(K),\qquad
P(A\cap K\cap U)\stackrel{?}{=}P(A)P(K)P(U).
\]

\textbf{Kendinize sorun}
\[
P(A\cup K\cup U)=P(A)+P(K)+P(U)-P(A\cap K)-P(A\cap U)-P(K\cap U)+P(A\cap K\cap U).
\]
}
\end{frame}

\begin{frame}[t,shrink=3]{Soru 4: Çözüm}
\scriptsize
\solutionsetup{ikili/tam bağımsızlık ve birleşim}{çarpım testi, dahil etme-dışlama}{önce bağımsızlık düzeyini netleştirip sonra birleşim formülüne geçmek gerekir.}
\solutionmistake{İkili bağımsızlık, otomatik olarak tam bağımsızlık anlamına gelmez; üçlü çarpım ayrıca kontrol edilir.}
\solutionstep{1}{İlk olarak her çift için çarpım koşulunu test ederek ikili bağımsızlığı doğruluyoruz.}
\solutionmath{
P(A\cap K)=P(A\cap U)=P(K\cap U)=\frac14
=\frac12\cdot\frac12.
}

\solutionstep{2}{Tam bağımsızlık için ayrıca \(P(A\cap K\cap U)=P(A)P(K)P(U)\) şartı sağlanmalıdır.}
\solutionmath{
P(A\cap K\cap U)=\frac18=P(A)P(K)P(U).
}

\solutionstep{3}{Birleşim olasılığını dahil etme-dışlama verir.}
\solutionmath{
P(A\cup K\cup U)=\frac12+\frac12+\frac12-\frac14-\frac14-\frac14+\frac18=\frac78.
}
\solutionfinal{Bu proje senaryosunda kararlar birbirini etkilemediği için hem ikili hem tam bağımsızlık sağlanır; en az bir projenin onay alma olasılığı da bu yüzden oldukça yüksektir.}
\end{frame}

\section{Koşullu Olasılık, Toplam Olasılık ve Bayes}
\subsection{Koşullu Yapılar}

\begin{frame}[t]{Soru 5: Sınıftan İki Öğrenci}
\footnotesize
Bir sınıfta 5 kadın ve 7 erkek öğrenci vardır. Rastgele iki öğrenci
\textbf{yerine koymadan} seçiliyor.

\textbf{İstenen}
\begin{itemize}
  \item[A.] İlk seçilen öğrenci kadın olduğuna göre, ikinci seçilenin erkek olma olasılığını bulunuz.
  \item[B.] Rastgele seçilen her iki öğrencinin de kadın olma olasılığını hesaplayınız.
  \item[C.] Öğrencilerden \textbf{en az birinin} kadın olduğu bilindiğine göre, her ikisinin de kadın olma olasılığını bulunuz.
\end{itemize}
\end{frame}

\begin{frame}[t]{Soru 5: Çözüm İçin Durak}
\small
\solutionghost{
\textbf{Kurulum}
\[
5\text{ kadın},\qquad 7\text{ erkek},\qquad \text{yerine koymadan 2 seçim}.
\]

\textbf{Kendinize sorun}
\begin{itemize}
  \item İlk bilginin verdiği yeni havuz kaç kişiden oluşuyor?
  \item “En az bir kadın” olayı doğrudan mı, tümleyenden mi daha rahat bulunur?
  \item Son adımda hangi koşullu olasılık tanımı kullanılacak?
\end{itemize}
}
\end{frame}

\begin{frame}[t,shrink=6]{Soru 5: Çözüm}
\scriptsize
\solutionsetup{yerine koymadan koşullu olasılık}{daralan örnek uzay, çarpım kuralı, koşullu olasılık tanımı}{ilk seçimin bilgisi sonraki olasılıkları değiştirir.}
\solutionmistake{Yeni bilgi geldikten sonra eski örnek uzayı kullanmayın; koşul havuzu küçültür.}
\solutionstep{1}{İlk bilgi havuzu daralttığı için ikinci seçimi yeni toplam üzerinden okuyoruz.}
\solutionmath{
P(\text{2. erkek}\mid \text{1. kadın})=\frac{7}{11}.
}

\solutionstep{2}{İkinci koşullu olasılık için pay ve paydayı ayrı ayrı kuruyoruz.}
\solutionmath{
P(\text{iki kadın})=\frac{5}{12}\cdot\frac{4}{11}=\frac{5}{33}.
}
\solutionmath{
P(\text{en az bir kadın})=1-P(\text{iki erkek})=1-\frac{7}{12}\cdot\frac{6}{11}=\frac{15}{22}.
}

\solutionstep{3}{Şimdi koşullu olasılık tanımını uyguluyoruz.}
\solutionmath{
P(\text{iki kadın}\mid \text{en az bir kadın})
=\frac{P(\text{iki kadın})}{P(\text{en az bir kadın})}
=\frac{5/33}{15/22}
=\frac{2}{9}.
}
\solutionfinal{Sınıf bağlamında “en az bir kadın” bilgisi olasılığı ciddi biçimde değiştirir; bu yüzden yerine koymadan seçimlerde verilen ek bilgi mutlaka yeni havuzla okunmalıdır.}
\end{frame}

\begin{frame}[t]{Soru 6: Kırmızı ve Mavi Toplar}
\footnotesize
Bir torbada 6 kırmızı ve 4 mavi top vardır. Art arda iki top çekiliyor;
yerine koyma yoktur.
İlk topun kırmızı gelmesi olayı \(K_1\), ikinci topun mavi gelmesi olayı \(M_2\) olsun.

\textbf{İstenen}
\begin{itemize}
  \item[A.] Ortak olasılığı bulunuz: \(P(K_1\cap M_2)\).
  \item[B.] Koşulsuz olasılığı bulunuz: \(P(M_2)\).
  \item[C.] Koşullu olasılığı bulunuz: \(P(M_2\mid K_1)\).
  \item[D.] \(K_1\) ve \(M_2\) olaylarının bağımsız olup olmadığını inceleyiniz.
\end{itemize}
\end{frame}

\begin{frame}[t]{Soru 6: Çözüm İçin Durak}
\small
\solutionghost{
\textbf{Kurulum}
\[
P(K_1)=\frac{6}{10},\qquad P(M_1)=\frac{4}{10}.
\]

\textbf{Kendinize sorun}
\begin{itemize}
  \item Ortak olasılık için hangi çarpım kuralı yazılmalı?
  \item \(P(M_2)\) doğrudan simetriyle mi, tam olasılıkla mı daha rahat bulunur?
  \item Bağımsızlık testinde hangi iki sayı karşılaştırılacak?
\end{itemize}
}
\end{frame}

\begin{frame}[t,shrink=8]{Soru 6: Çözüm}
\scriptsize
\solutionsetup{koşullu/koşulsuz olasılık ve bağımsızlık}{çarpım kuralı, tam olasılık, bağımsızlık testi}{aynı torba örneğinde ortak, koşulsuz ve koşullu olasılıklar farklı niceliklerdir.}
\solutionmistake{\(P(A\cap B)\), \(P(A)\) ve \(P(A\mid B)\) birbirinin yerine yazılamaz; önce hangi niceliğin istendiğini ayırın.}
\solutionstep{A}{Ortak olasılığı çarpım kuralıyla yazıyoruz.}
\solutionmath{
P(K_1\cap M_2)=P(K_1)P(M_2\mid K_1)=\frac{6}{10}\cdot\frac{4}{9}=\frac{4}{15}.
}

\solutionstep{B}{İkinci topun mavi gelmesi olasılığı koşulsuz bir olasılıktır; hem doğrudan hem tam olasılıkla bulunabilir.}
\solutionmath{
P(M_2)=\frac{4}{10}=\frac{2}{5}
=\frac49\cdot\frac{6}{10}+\frac39\cdot\frac{4}{10}.
}

\solutionstep{C}{İlk topun kırmızı geldiği bilgisi verildiğinde havuz 9 topa düşer ve bunların 4'ü mavidir.}
\solutionmath{
P(M_2\mid K_1)=\frac{4}{9}.
}

\solutionstep{D}{Bağımsızlık için \(P(K_1\cap M_2)=P(K_1)P(M_2)\) şartı sağlanmalıdır.}
\solutionmath{
P(K_1\cap M_2)=\frac{4}{15},\qquad P(K_1)P(M_2)=\frac{6}{10}\cdot\frac{2}{5}=\frac{6}{25}.
}
\solutionfinal{Torba bağlamında ilk çekiliş ikinciyi etkilediği için ortak, koşulsuz ve koşullu olasılıklar farklı çıkar; bu da olayların bağımsız olmadığını gösterir.}
\end{frame}

\begin{frame}[t]{Tablo Temelli Örnek: Sertifika Alt Uzayında Oran}
\footnotesize
\begin{columns}[T,onlytextwidth]
    \begin{column}{0.56\textwidth}
        Bir elektronik üretim tesisindeki teknisyenlerin ekip ve kalibrasyon sertifikası dağılımı:
        
        \begin{table}[]
            \centering
            \resizebox{0.88\textwidth}{!}{%
            \begin{tabular}{@{}lccc@{}}
                \toprule
                & \textbf{Sertifikalı (K)} & \textbf{Sertifikasız} & \textbf{Toplam} \\ \midrule
                \textbf{Otomasyon Ekibi (O)} & 150 & 30 & 180 \\
                \textbf{Bakım Ekibi (B)} & 80 & 140 & 220 \\ \midrule
                \textbf{Toplam} & \textbf{230} & \textbf{170} & \textbf{400} \\ \bottomrule
            \end{tabular}%
            }
        \end{table}
    
        \small
        \textbf{Soru:} Rastgele seçilen bir teknisyenin \textbf{sertifikalı (K)} olduğu biliniyorsa, bu kişinin \textbf{otomasyon ekibinde (O)} olma olasılığı nedir?
    \end{column}
    
    \begin{column}{0.42\textwidth}
        \onslide<2->{
            \scriptsize
            \textbf{Çözüm (İndirgenmiş Uzay):}
            
            Sadece ``Sertifikalı'' sütununa bakılır. Toplam \(230\) sertifikalı teknisyen vardır. Bunların \(150\)'si otomasyon ekibindedir.
            \[
            P(O\mid K)=\frac{150}{230}=\frac{15}{23}\approx 0.652
            \]
            
            \vspace{0.2cm}
            \textbf{Formül ile:}
            \[
            P(K)=\frac{230}{400},\qquad P(O\cap K)=\frac{150}{400}
            \]
            \[
            P(O\mid K)=\frac{P(O\cap K)}{P(K)}
            =\frac{150/400}{230/400}
            =\frac{150}{230}
            \]
            
            Sertifikalı bir teknisyen seçildiğinde otomasyon ekibinden gelme olasılığı yaklaşık \(\%65.2\)'dir.
        }
    \end{column}
\end{columns}
\end{frame}

\begin{frame}[t]{Soru 7: Üç Makine ve Bayes}
\footnotesize
Bir ürünün 1., 2. ve 3. makineden gelme olayları sırasıyla \(M_1\), \(M_2\), \(M_3\),
ürünün hasarlı olması olayı ise \(H\) olsun.
\[
P(M_1)=0.25,\quad P(M_2)=0.50,\quad P(M_3)=0.25,
\]
\[
P(H\mid M_1)=0.04,\quad P(H\mid M_2)=0.02,\quad P(H\mid M_3)=0.05.
\]

\textbf{İstenen}
\begin{itemize}
  \item[A.] \(P(H)\) bulunuz.
  \item[B.] \(P(M_3\mid H)\) bulunuz.
  \item[C.] \(P(M_2\mid H')\) bulunuz.
\end{itemize}
\end{frame}

\begin{frame}[t]{Soru 7: Çözüm İçin Durak}
\small
\solutionghost{
\textbf{Kurulum}
\[
P(M_1)=0.25,\quad P(M_2)=0.50,\quad P(M_3)=0.25.
\]
\[
P(H\mid M_1)=0.04,\quad P(H\mid M_2)=0.02,\quad P(H\mid M_3)=0.05.
\]

\textbf{İlk hangi kural?}
\[
P(H)=\sum_{i=1}^{3} P(H\mid M_i)P(M_i).
\]

\textbf{Kendinize sorun}
\begin{itemize}
  \item Bayes paydasında önce hangi toplam olasılık bulunmalı?
  \item \(H'\) için hangi koşullu olasılıklar kullanılacak?
\end{itemize}
}
\end{frame}

\begin{frame}[t,shrink=7]{Soru 7: Çözüm}
\footnotesize
\solutionsetup{toplam olasılık ve Bayes}{ağırlıklı toplam, Bayes normalizasyonu}{gözlenen kusurdan kaynağa dönmek için önce toplam kusur oranı gerekir.}
\solutionmistake{Bayes'te yön karıştırmayın; \(P(M_3\mid H)\) için pay kısmında \(P(H\mid M_3)\) gerekir.}
\solutionstep{1}{Önce toplam hasarlı olasılığını buluyoruz; çünkü Bayes formülündeki payda bu değerdir.}
\solutionmath{
P(H)=\sum_{i=1}^{3}P(H\mid M_i)P(M_i)
=0.25(0.04)+0.50(0.02)+0.25(0.05)
=0.0325.
}

\solutionstep{2}{Hasarlı ürün gözlendiğinde Bayes kuralını uyguluyoruz.}
\solutionmath{
P(M_3\mid H)=\frac{P(H\mid M_3)P(M_3)}{P(H)}
=\frac{0.05\cdot 0.25}{0.0325}
=\frac{5}{13}\approx 0.3846.
}

\solutionstep{3}{Hasarsız ürün için yine Bayes kuralını bu kez \(H'\) üzerinden yazıyoruz.}
\solutionmath{
P(M_2\mid H')=\frac{P(H'\mid M_2)P(M_2)}{P(H')}
=\frac{0.98\cdot 0.50}{1-0.0325}
=\frac{196}{387}\approx 0.5065.
}
\solutionfinal{Hasarlı ürün görüldüğünde üçüncü makinenin payı yaklaşık \(\%38\) olur; hasarsız üründe ise ikinci makine yaklaşık \(\%51\) ile en olası kaynak görünür.}
\end{frame}

\begin{frame}[t]{Soru 8: Hastalık Testi}
\footnotesize
Bir tarama testinde kişinin hasta olma olayı \(H\), test sonucunun pozitif çıkması olayı \(+\) ile gösterilsin.
\[
P(H)=0.03,\qquad P(+\mid H)=0.92,\qquad P(+\mid H')=0.07.
\]

\textbf{İstenen}
\begin{itemize}
  \item[A.] \(P(+)\) bulunuz.
  \item[B.] \(P(H\mid +)\) bulunuz.
  \item[C.] \(P(H\mid -)\) bulunuz.
  \item[D.] Yorum yapınız.
\end{itemize}
\end{frame}

\begin{frame}[t]{Soru 8: Çözüm İçin Durak}
\small
\solutionghost{
\textbf{Kurulum}
\[
P(H)=0.03,\qquad P(+\mid H)=0.92,\qquad P(+\mid H')=0.07.
\]

\textbf{İlk hangi kural?}
\[
P(+)=P(+\mid H)P(H)+P(+\mid H')P(H').
\]

\textbf{Kendinize sorun}
\begin{itemize}
  \item Bayes paydasında pozitif testin toplam olasılığı nasıl bulunur?
  \item Negatif sonuç için \(P(-\mid H)\) ve \(P(-)\) nasıl yazılır?
\end{itemize}
}
\end{frame}

\begin{frame}[t,shrink=6]{Soru 8: Çözüm}
\footnotesize
\solutionsetup{Bayes ve tanı testi yorumu}{toplam olasılık, Bayes, tamamlayıcı olay}{test sonucu ile gerçek durum ters yönde okunacağı için önce toplam pozitiflik gerekir.}
\solutionmistake{Duyarlılık \(P(+\mid H)\) ile pozitif çıkanın hasta olma olasılığı \(P(H\mid +)\) aynı değildir.}
\solutionstep{1}{Pozitif test sonucu, hasta ve sağlıklı dallarının toplamından gelir; önce bu paydayı kuruyoruz.}
\solutionmath{
P(+)=P(+\mid H)P(H)+P(+\mid H')P(H')
=0.92(0.03)+0.07(0.97)
=0.0955.
}

\solutionstep{2}{Pozitif çıkan kişinin gerçekten hasta olma olasılığını Bayes ile buluyoruz.}
\solutionmath{
P(H\mid +)=\frac{P(+\mid H)P(H)}{P(+)}
=\frac{0.92\cdot 0.03}{0.0955}
=\frac{276}{955}\approx 0.289.
}

\solutionstep{3}{Negatif sonuç için aynı mantığı tamamlayıcı olayla kuruyoruz.}
\solutionmath{
P(H\mid -)=\frac{P(-\mid H)P(H)}{P(-)}
=\frac{0.08\cdot 0.03}{1-0.0955}
\approx 0.00265.
}
\solutionfinal{Bu testte pozitif sonuç alan kişinin gerçekten hasta olma olasılığı yaklaşık \(\%29\)'dur; yani düşük prevalans, yüksek duyarlılığa rağmen pozitif sonucun güvenini sınırlamaktadır.}
\end{frame}

\begin{frame}[t]{Neden Şimdi Rassal Değişken?}
\small
\begin{itemize}
  \item Buraya kadar olayların olasılığını hesapladık; şimdi her sonucu bir sayıya dönüştüren \(X\) ile çalışacağız.
  \item Yeni hedef, sadece ``olay oldu mu?'' değil, ``hangi değer ne kadar olası?'' sorusunu tek bir dağılım nesnesiyle okumaktır.
  \item Bu bölümde PMF, CDF ve PDF; rassal değişkenin ayrık mı sürekli mi olduğunu netleştirmek için kullanılacak.
\end{itemize}

\begin{infoblock}{Neden bu geçiş önemli?}
Bayes ve koşullu olasılık bize olaylar arasında yön seçmeyi öğretti; şimdi aynı düşünceyi bir dağılımı okumaya ve birikimli olasılığı yorumlamaya taşıyoruz.
\end{infoblock}
\end{frame}

\section{Rassal Değişkenler ve Dağılımlar}
\subsection{PMF ve CDF}

\begin{frame}[t]{Soru 9: 8 Noktalı Örnek Uzaydan \(X\)}
\footnotesize
\[
S=\{1,2,3,4,5,6,7,8\}
\]
ve tüm temel sonuçlar eş olasılıklıdır.

\[
X(\omega)=
\begin{cases}
0, & \omega\in\{1,2\},\\
1, & \omega\in\{3,4,5\},\\
2, & \omega\in\{6,7\},\\
4, & \omega=8.
\end{cases}
\]

\textbf{İstenen}
\begin{itemize}
  \item[A.] PMF'yi bulunuz.
  \item[B.] \(F(x)\) fonksiyonunu parçalı yazınız.
  \item[C.] \(P(1\le X<4)\) ve \(P(X\neq 2)\) bulunuz.
\end{itemize}
\end{frame}

\begin{frame}[t]{Soru 9: Çözüm İçin Durak}
\small
\solutionghost{
\textbf{Kurulum}
\[
X(\omega)=
\begin{cases}
0, & \omega\in\{1,2\},\\
1, & \omega\in\{3,4,5\},\\
2, & \omega\in\{6,7\},\\
4, & \omega=8.
\end{cases}
\]

\textbf{Kendinize sorun}
\[
\begin{array}{c|cccc}
x & 0 & 1 & 2 & 4\\ \hline
f(x) & ? & ? & ? & ? \\
F(x) & ? & ? & ? & ?
\end{array}
\]
\begin{itemize}
  \item Her \(x\) değerini kaç temel sonuç üretiyor?
  \item CDF'yi soldan sağa hangi sırayla biriktireceksiniz?
\end{itemize}
}
\end{frame}

\begin{frame}[t,shrink=6]{Soru 9: Çözüm}
\scriptsize
\solutionsetup{ayrık rassal değişken, PMF ve CDF}{eş olasılıklı sayma, birikimli toplama}{önce tekil kütleler bulunmadan CDF kurulamaz.}
\solutionmistake{CDF değeri tek noktadaki kütle değil, o noktaya kadar birikmiş olasılıktır.}
\solutionstep{1}{Önce her \(x\) değerini üreten temel sonuç sayısını sayarak PMF'yi yazıyoruz.}
\solutionmath{
P(X=0)=\frac28=\frac14,\qquad
P(X=1)=\frac38,\qquad
P(X=2)=\frac28=\frac14,\qquad
P(X=4)=\frac18.
}

\solutionstep{2}{Sonra bu kütleleri soldan sağa toplayarak CDF'yi açıkça kuruyoruz.}
\solutionmath{
F(x)=
\begin{cases}
0, & x<0,\\
\frac14, & 0\le x<1,\\
\frac58, & 1\le x<2,\\
\frac78, & 2\le x<4,\\
1, & x\ge 4.
\end{cases}
}

\solutionstep{3}{İstenen aralık olasılıklarını PMF üzerinden topluyoruz.}
\solutionmath{
P(1\le X<4)=P(X=1)+P(X=2)=\frac38+\frac14=\frac58,\qquad
P(X\neq 2)=1-\frac14=\frac34.
}
\solutionfinal{Bu örnek uzayda rassal değişken yalnızca dört değer alabildiği için dağılım önce sayma ile, sonra soldan birikimli toplama ile okunur.}
\end{frame}

\begin{frame}[t]{Soru 10: \(f(x)=c(x+1)\)}
\footnotesize
Ayrık bir rassal değişken \(X\) için olasılık kütle fonksiyonu aşağıdaki gibi veriliyor:
\[
f(x)=c(x+1),\qquad x=0,1,2,3,
\]
ve diğer durumda \(0\).

\textbf{İstenen}
\begin{itemize}
  \item[A.] \(c\) sabitini bulunuz.
  \item[B.] \(P(X\le 2)\) ve \(P(1<X\le 3)\) bulunuz.
  \item[C.] \(F(x)\) fonksiyonunu yazınız.
\end{itemize}
\end{frame}

\begin{frame}[t]{Soru 10: Çözüm İçin Durak}
\small
\solutionghost{
\textbf{Kurulum}
\[
f(x)=c(x+1),\qquad x=0,1,2,3.
\]

\textbf{İlk hangi adım?}
\[
\sum_{x=0}^{3} f(x)=1.
\]

\textbf{Kendinize sorun}
\[
\begin{array}{c|cccc}
x & 0 & 1 & 2 & 3\\ \hline
f(x) & ? & ? & ? & ? \\
F(x) & ? & ? & ? & ?
\end{array}
\]
\begin{itemize}
  \item Önce \(c\) bulunmadan hangi olasılıklar hesaplanamaz?
  \item CDF tablosu için hangi sırayla biriktirme yapacaksınız?
\end{itemize}
}
\end{frame}

\begin{frame}[t,shrink=6]{Soru 10: Çözüm}
\scriptsize
\solutionsetup{parametreli PMF ve CDF}{normlama, tekil olasılıkları toplama, birikimli toplama}{bilinmeyen sabit bulunmadan ne olasılık ne CDF yazılabilir.}
\solutionmistake{Önce normlama sabitini bulmadan CDF parçalarını yazmaya başlamayın.}
\solutionstep{1}{İlk olarak normlama koşulunu kullanıp sabiti belirliyoruz; çünkü dağılım önce geçerli olmalıdır.}
\solutionmath{
c[(0+1)+(1+1)+(2+1)+(3+1)]=10c=1
\Longrightarrow c=\frac{1}{10}.
}

\solutionstep{2}{İstenen olasılıkları ilgili kütleleri toplayarak hesaplıyoruz.}
\solutionmath{
P(X\le 2)=\frac{1+2+3}{10}=\frac35,\qquad
P(1<X\le 3)=\frac{3+4}{10}=\frac{7}{10}.
}

\solutionstep{3}{Son olarak PMF'den soldan toplama yaparak CDF'yi açıkça yazıyoruz.}
\solutionmath{
F(x)=
\begin{cases}
0, & x<0,\\
\frac{1}{10}, & 0\le x<1,\\
\frac{3}{10}, & 1\le x<2,\\
\frac{6}{10}, & 2\le x<3,\\
1, & x\ge 3.
\end{cases}
}
\solutionfinal{Bu dağılımda kütle \(x\) büyüdükçe arttığı için CDF de küçük değerlerde yavaş, büyük değerlere yaklaştıkça daha hızlı yükselir.}
\end{frame}

\begin{frame}[t]{Soru 11: CDF'den PMF'ye Geçiş}
\footnotesize
\[
F(x)=
\begin{cases}
0, & x<0,\\
0.2, & 0\le x<1,\\
0.5, & 1\le x<3,\\
0.9, & 3\le x<5,\\
1, & x\ge 5.
\end{cases}
\]

\textbf{İstenen}
\begin{itemize}
  \item[A.] \(X\)'in değerlerini ve olasılıklarını bulunuz.
  \item[B.] \(P(X=3)\), \(P(X>1)\), \(P(0\le X<5)\) bulunuz.
  \item[C.] \(X\) ayrık mı, sürekli mi?
\end{itemize}
\end{frame}

\begin{frame}[t]{Soru 11: Çözüm İçin Durak}
\small
\solutionghost{
\textbf{Kurulum}
\[
F(x)=
\begin{cases}
0, & x<0,\\
0.2, & 0\le x<1,\\
0.5, & 1\le x<3,\\
0.9, & 3\le x<5,\\
1, & x\ge 5.
\end{cases}
\]

\textbf{Kendinize sorun}
\[
\begin{array}{c|cccc}
x & 0 & 1 & 3 & 5\\ \hline
\text{CDF sıçraması} & ? & ? & ? & ? \\
P(X=x) & ? & ? & ? & ?
\end{array}
\]
\begin{itemize}
  \item Her sıçrama hangi noktanın PMF değeridir?
  \item \(P(X>1)\) için hangi destek noktaları toplanmalıdır?
\end{itemize}
}
\end{frame}

\begin{frame}[t,shrink=3]{Soru 11: Çözüm}
\scriptsize
\solutionsetup{CDF'den PMF'ye geçiş}{sıçrama büyüklüğü, destek okuma, kütle toplama}{veri CDF biçiminde verildiği için önce sıçramaları kütleye çevirmek gerekir.}
\solutionmistake{PMF'yi CDF yüksekliğinden değil, ardışık CDF seviyeleri arasındaki sıçrama farkından çıkarın.}
\solutionstep{1}{Ayrık değişkende her sıçrama miktarı ilgili noktadaki PMF değeridir; önce bunu tabloya döküyoruz.}
\solutionmath{
P(X=0)=0.2-0=0.2,\quad
P(X=1)=0.5-0.2=0.3,\quad
P(X=3)=0.9-0.5=0.4,\quad
P(X=5)=1-0.9=0.1.
}
\solutionmath{
\begin{array}{c|cccc}
x & 0 & 1 & 3 & 5\\ \hline
P(X=x) & 0.2 & 0.3 & 0.4 & 0.1
\end{array}
}

\solutionstep{2}{Böylece destek kümesini ve istenen olasılıkları doğrudan yazabiliyoruz.}
\solutionmath{
X\in\{0,1,3,5\},\qquad
P(X=3)=0.4.
}
\solutionmath{
P(X>1)=P(X=3)+P(X=5)=0.4+0.1=0.5,\qquad
P(0\le X<5)=0.2+0.3+0.4=0.9.
}
\solutionfinal{Burada dağılım yalnızca sıçrama noktalarında kütle taşıdığı için \(X\) ayrık bir değişkendir; CDF grafiği de bu nedenle basamaklı görünür.}
\end{frame}

\subsection{PDF}

\begin{frame}[t]{Soru 12: \(k x^2\) Yoğunluğu}
\footnotesize
Sürekli bir rassal değişken \(X\) için yoğunluk fonksiyonu aşağıdaki gibi veriliyor:
\[
f(x)=
\begin{cases}
kx^2, & 0<x<2,\\
0, & \text{diğer}.
\end{cases}
\]

\textbf{İstenen}
\begin{itemize}
  \item[A.] \(k\) bulunuz.
  \item[B.] \(P(0.5<X<1.5)\) hesaplayınız.
  \item[C.] \(F(x)\) fonksiyonunu yazınız.
  \item[D.] \(P(X=1)\) değerini açıklayınız.
\end{itemize}
\end{frame}

\begin{frame}[t]{Soru 12: Çözüm İçin Durak}
\small
\solutionghost{
\textbf{Kurulum}
\[
f(x)=
\begin{cases}
kx^2, & 0<x<2,\\
0, & \text{diğer}.
\end{cases}
\]

\textbf{İlk hangi adım?}
\[
\int_0^2 kx^2\,dx=1.
\]

\textbf{Kendinize sorun}
\begin{itemize}
  \item Aralık olasılığı için hangi integral kurulmalı?
  \item CDF yazarken \(F(x)=\int_{-\infty}^{x} f(t)\,dt\) hangi üç bölgeye ayrılır?
\end{itemize}
}
\end{frame}

\begin{frame}[t,shrink=4]{Soru 12: Çözüm}
\scriptsize
\solutionsetup{PDF, integral olasılığı ve CDF}{normlama integrali, alan hesabı, soldan integral}{sürekli dağılımda tüm sorular yoğunluğun integrali üzerinden çözülür.}
\solutionmistake{Sürekli dağılımda \(f(1)\) bir olasılık değildir; \(P(X=1)=0\) kalır.}
\solutionstep{1}{Önce normlama integralini kuruyoruz; çünkü \(k\) bilinmeden yoğunluk tamamlanmaz.}
\solutionmath{
1=\int_0^2 kx^2\,dx
=k\left[\frac{x^3}{3}\right]_0^2
=\frac{8k}{3}
\Longrightarrow k=\frac{3}{8}.
}

\solutionstep{2}{İstenen aralık olasılığını integral ile hesaplıyoruz.}
\solutionmath{
P(0.5<X<1.5)=\int_{0.5}^{1.5}\frac{3}{8}x^2\,dx
=\left[\frac{x^3}{8}\right]_{0.5}^{1.5}
=\frac{13}{32}.
}

\solutionstep{3}{CDF'yi ve tek nokta olasılığını yazıyoruz.}
\solutionmath{
F(x)=
\begin{cases}
0, & x\le 0,\\
\frac{x^3}{8}, & 0<x<2,\\
1, & x\ge 2.
\end{cases}
}
\solutionfinal{Bu yoğunlukta olasılık alanla ölçüldüğü için tek bir noktanın olasılığı sıfırdır; yani \(X=1\) tam değeri değil, aralıklar anlam taşır.}
\end{frame}

\begin{frame}[t]{Soru 13: \(c(2-x)\) Yoğunluğu}
\footnotesize
Sürekli bir rassal değişken \(X\) için yoğunluk fonksiyonu aşağıdaki gibi veriliyor:
\[
f(x)=
\begin{cases}
c(2-x), & 0<x<2,\\
0, & \text{diğer}.
\end{cases}
\]

\textbf{İstenen}
\begin{itemize}
  \item[A.] \(c\) bulunuz.
  \item[B.] \(P(X<1)\) ve \(P(X\ge 1.5)\) bulunuz.
\end{itemize}
\end{frame}

\begin{frame}[t]{Soru 13: Çözüm İçin Durak}
\small
\solutionghost{
\textbf{Kurulum}
\[
f(x)=
\begin{cases}
c(2-x), & 0<x<2,\\
0, & \text{diğer}.
\end{cases}
\]

\textbf{Kendinize sorun}
\begin{itemize}
  \item Önce \(c\) için hangi normlama integrali yazılmalı?
  \item İstenen iki olasılık hangi alt aralıkların alanıdır?
\end{itemize}
}
\end{frame}

\begin{frame}[t]{Soru 13: Çözüm}
\footnotesize
\solutionsetup{parametreli PDF ve alan olasılığı}{normlama integrali, alt aralık alanı}{yoğunluk önce geçerli hale getirilmeli, sonra alanlar hesaplanmalıdır.}
\solutionmistake{Aralık olasılığını tek bir \(f(x)\) değeriyle değil, ilgili bölgenin alanıyla hesaplayın.}
\solutionstep{1}{İlk olarak toplam alanı \(1\)'e eşitleyip \(c\) sabitini belirliyoruz.}
\solutionmath{
1=\int_0^2 c(2-x)\,dx
=c\left[2x-\frac{x^2}{2}\right]_0^2
=2c
\Longrightarrow c=\frac12.
}

\solutionstep{2}{İstenen olasılıkları doğrudan integral ile hesaplıyoruz.}
\solutionmath{
P(X<1)=\int_0^1 \frac{2-x}{2}\,dx=\frac34,\qquad
P(X\ge 1.5)=\int_{1.5}^{2}\frac{2-x}{2}\,dx=\frac{1}{16}.
}
\solutionfinal{Bu yoğunluk soldan sağa azaldığı için küçük \(x\) değerleri daha olasıdır; bu yüzden \(X<1\) olasılığı sağ kuyruktaki \(X\ge 1.5\) olasılığından çok daha büyüktür.}
\end{frame}

\begin{frame}[t]{Soru 14: Geçerli PDF mi?}
\footnotesize
Bir sürekli rassal değişken \(X\) için önerilen yoğunluk fonksiyonu aşağıdadır:
\[
f(x)=
\begin{cases}
\dfrac{2(x+2)}{5}, & 0<x<1,\\
0, & \text{diğer}.
\end{cases}
\]

\textbf{İstenen}
\begin{itemize}
  \item[A.] Geçerli bir yoğunluk olduğunu gösteriniz.
  \item[B.] \(P(0.25<X<0.75)\) bulunuz.
  \item[C.] \(F(x)\) fonksiyonunu yazınız.
\end{itemize}
\end{frame}

\begin{frame}[t]{Soru 14: Çözüm İçin Durak}
\small
\solutionghost{
\textbf{Kurulum}
\[
f(x)=
\begin{cases}
\dfrac{2(x+2)}{5}, & 0<x<1,\\
0, & \text{diğer}.
\end{cases}
\]

\textbf{Kontrol listesi}
\begin{itemize}
  \item Önce negatif olmama koşulunu inceleyin.
  \item Sonra toplam alanın \(1\) olup olmadığını test edin.
  \item Ardından aralık olasılığı ve CDF için integral yazın.
\end{itemize}
}
\end{frame}

\begin{frame}[t,shrink=4]{Soru 14: Çözüm}
\scriptsize
\solutionsetup{geçerli PDF testi ve CDF kurma}{işaret kontrolü, normlama integrali, soldan integral}{bir fonksiyonun yoğunluk sayılabilmesi için önce geçerlilik şartları doğrulanmalıdır.}
\solutionmistake{Toplam alan \(1\) olsa bile negatif kalan bir fonksiyon PDF olamaz; işaret kontrolünü atlamayın.}
\solutionstep{1}{İlk işimiz PDF koşullarını test etmektir; çünkü ancak geçerli yoğunluk üzerinde olasılık konuşabiliriz.}
\solutionmath{
\frac{2(x+2)}{5}>0 \quad (0<x<1),\qquad
\int_0^1 \frac{2(x+2)}{5}\,dx=1.
}

\solutionstep{2}{İstenen aralık olasılığını integral ile buluyoruz.}
\solutionmath{
P(0.25<X<0.75)=\int_{0.25}^{0.75}\frac{2(x+2)}{5}\,dx=\frac12.
}

\solutionstep{3}{CDF'yi tanım gereği soldan integralle kuruyoruz.}
\solutionmath{
F(x)=
\begin{cases}
0, & x\le 0,\\
\int_0^x \frac{2(t+2)}{5}\,dt=\frac{x^2}{5}+\frac{4x}{5}, & 0<x<1,\\
1, & x\ge 1.
\end{cases}
}
\solutionfinal{Bu fonksiyon \((0,1)\) aralığında her yerde pozitif ve toplam alanı \(1\) olduğundan, gerçekten de bir olasılık yoğunluğu gibi davranır.}
\end{frame}

\begin{frame}[t]{Neden Şimdi Beklenen Değer ve Varyans?}
\small
\begin{itemize}
  \item PMF ve PDF bize dağılımın tamamını verdi; şimdi bu dağılımı birkaç özet sayı ile yorumlamayı öğreneceğiz.
  \item Beklenen değer merkezin nerede toplandığını, varyans ise saçılmanın ne kadar güçlü olduğunu söyler.
  \item Bu bölümde aynı dağılımdan moment, dönüşüm ve karşılaştırma soruları çıkaracağız.
\end{itemize}

\begin{infoblock}{Yeni karar sorusu}
Bir dağılımı yazabiliyorsanız artık ``ortalama ne?'', ``yayılım ne kadar?'' ve ``dönüşüm bu iki niceliği nasıl değiştirir?'' sorularını sorabilirsiniz.
\end{infoblock}
\end{frame}

\section{Beklenen Değer ve Varyans}
\subsection{Ayrık Beklenen Değer}

\begin{frame}[t]{Soru 15: \(-1,0,2\) Dağılımı}
\footnotesize
Ayrık bir rassal değişken \(X\) yalnızca \(-1\), \(0\) ve \(2\) değerlerini alıyor.
\[
P(X=-1)=0.2,\qquad P(X=0)=0.3,\qquad P(X=2)=0.5.
\]

\textbf{İstenen}
\begin{itemize}
  \item[A.] \(E(X)\) bulunuz.
  \item[B.] \(E(X^2)\) ve \(E((2X+1)^2)\) bulunuz.
  \item[C.] \(g(X)=3X-4\) için \(E(g(X))\)'i hem doğrudan hem özellik ile doğrulayınız.
\end{itemize}

\begin{infoblock}{Hatırlatma}
Ayrık durumda
\[
E(X)=\sum_x x\,f(x),\qquad E(g(X))=\sum_x g(x)\,f(x).
\]
\end{infoblock}
\end{frame}

\begin{frame}[t]{Soru 15: Çözüm İçin Durak}
\small
\solutionghost{
\textbf{Kurulum}
\[
\begin{array}{c|ccc}
X & -1 & 0 & 2\\ \hline
P(X=x) & 0.2 & 0.3 & 0.5
\end{array}
\]

\textbf{Kendinize sorun}
\begin{itemize}
  \item Hangi satır \(E(X)\), hangisi \(E(X^2)\) için gerekir?
  \item \(2X+1\) ve \(3X-4\) dönüşümlerini destek noktalarında tek tek yazdınız mı?
  \item Son maddede doğrudan hesap ile lineerlik özelliği aynı sonucu veriyor mu?
\end{itemize}
}
\end{frame}

\begin{frame}[t,shrink=8]{Soru 15: Çözüm}
\scriptsize
\solutionsetup{ayrık beklenen değer ve dönüşüm}{beklenen değer toplamı, tablo üzerinden dönüşüm}{destek noktaları az olduğu için en güvenli yol tabloyla doğrudan hesaplamadır.}
\solutionmistake{\(E(g(X))\) genel olarak \(g(E(X))\) değildir; dönüşümü önce destek noktalarında üretin.}
\solutionstep{1}{Önce beklenen değer ve ikinci momenti aynı tablo üzerinden doğrudan topluyoruz.}
\[
\begin{array}{c|ccc}
X & -1 & 0 & 2\\ \hline
X^2 & 1 & 0 & 4\\
(2X+1)^2 & 1 & 1 & 25\\
3X-4 & -7 & -4 & 2
\end{array}
\]

\solutionmath{
E(X)=(-1)(0.2)+0(0.3)+2(0.5)=0.8,\qquad
E(X^2)=(-1)^2(0.2)+2^2(0.5)=2.2.
}

\solutionstep{2}{Dönüşmüş kare için her destek noktasını yerine yazıyoruz.}
\solutionmath{
E((2X+1)^2)=1(0.2)+1(0.3)+25(0.5)=13.
}

\solutionstep{3}{\(g(X)=3X-4\) için doğrudan ve özellik yolunu karşılaştırıyoruz.}
\solutionmath{
E(3X-4)=(-7)(0.2)+(-4)(0.3)+2(0.5)=-1.6
=3E(X)-4.
}
\solutionfinal{Bu dağılımda doğrusal dönüşüm hem doğrudan hem özellik yoluyla aynı sonucu verdiği için, lineerlik kuralının pratikte nasıl çalıştığı açıkça görülür.}
\end{frame}

\begin{frame}[t]{Soru 16: \(x=\{0,1,3,5\}\) Dağılımı}
\footnotesize
Ayrık bir rassal değişken \(X\)'in kütle fonksiyonu tabloyla veriliyor:
\[
\begin{array}{c|cccc}
x & 0 & 1 & 3 & 5\\ \hline
f(x) & 0.1 & 0.4 & 0.3 & 0.2
\end{array}
\]

\textbf{İstenen}
\begin{itemize}
  \item[A.] \(E(X)\) bulunuz.
  \item[B.] \(\mathrm{Var}(X)\) ve \(\sigma_X\) bulunuz.
  \item[C.] \(Y=2X-1\) için \(E(Y)\) ve \(\mathrm{Var}(Y)\) bulunuz.
\end{itemize}

\begin{infoblock}{Hatırlatma}
Ayrık durumda
\[
E(X)=\sum_x x\,f(x),\qquad E(Y)=\sum_x y\,f_Y(y).
\]
\[
\mathrm{Var}(X)=\sum_x x^2 f(x)-[E(X)]^2=E(X^2)-[E(X)]^2.
\]
\end{infoblock}
\end{frame}

\begin{frame}[t]{Soru 16: Çözüm İçin Durak}
\small
\solutionghost{
\textbf{Kurulum}
\[
\begin{array}{c|cccc}
x & 0 & 1 & 3 & 5\\ \hline
f(x) & 0.1 & 0.4 & 0.3 & 0.2
\end{array}
\]

\textbf{Kendinize sorun}
\begin{itemize}
  \item Önce \(E(X)\) ve \(E(X^2)\) için hangi iki toplam yazılmalı?
  \item Varyans için doğrudan hangi fark kullanılacak?
  \item \(Y=2X-1\) için ortalama ve varyans nasıl ölçeklenir?
\end{itemize}
}
\end{frame}

\begin{frame}[t,shrink=3]{Soru 16: Çözüm}
\scriptsize
\solutionsetup{ayrık beklenen değer, varyans ve dönüşüm}{moment toplamları ve dönüşüm kuralları}{varyans için önce birinci ve ikinci moment gerekir.}
\solutionmistake{\(\mathrm{Var}(2X-1)\), \(2\mathrm{Var}(X)-1\) değildir; sabit terim varyansı değiştirmez, katsayı kareyle gelir.}
\solutionstep{1}{İlk adımda birinci ve ikinci momentleri hesaplıyoruz; çünkü varyans doğrudan bunlardan çıkacak.}
\solutionmath{
E(X)=0(0.1)+1(0.4)+3(0.3)+5(0.2)=2.3,\qquad
E(X^2)=1^2(0.4)+3^2(0.3)+5^2(0.2)=8.1.
}

\solutionstep{2}{Buradan varyans ve standart sapma gelir.}
\solutionmath{
\mathrm{Var}(X)=8.1-(2.3)^2=2.81,\qquad
\sigma_X=\sqrt{2.81}\approx 1.676.
}

\solutionstep{3}{\(Y=2X-1\) için dönüşüm özelliklerini uyguluyoruz.}
\solutionmath{
E(Y)=2E(X)-1=3.6,\qquad
\mathrm{Var}(Y)=4\,\mathrm{Var}(X)=11.24.
}
\solutionfinal{Bu tabloda dönüşüm sonrası ortalama doğrusal değişirken varyansın daha hızlı büyümesi, ölçeklemenin yayılımı nasıl etkilediğini gösterir.}
\end{frame}

\subsection{Sürekli Beklenen Değer}

\begin{frame}[t]{Soru 17: \(f(x)=x/2\)}
\footnotesize
Sürekli bir rassal değişken \(X\) için yoğunluk fonksiyonu aşağıdaki gibidir:
\[
f(x)=
\begin{cases}
\dfrac{x}{2}, & 0<x<2,\\
0, & \text{diğer}.
\end{cases}
\]

\textbf{İstenen}
\begin{itemize}
  \item[A.] \(E(X)\), \(E(X^2)\) bulunuz.
  \item[B.] \(\mathrm{Var}(X)\) ve \(\sigma_X\) hesaplayınız.
  \item[C.] \(g(X)=4X+3\) için \(E(g(X))\) bulunuz.
\end{itemize}

\begin{infoblock}{Hatırlatma}
Sürekli durumda
\[
E(X)=\int_{-\infty}^{\infty} x\,f(x)\,dx,\qquad
E(g(X))=\int_{-\infty}^{\infty} g(x)\,f(x)\,dx.
\]
\[
\mathrm{Var}(X)=\int_{-\infty}^{\infty} x^2 f(x)\,dx-[E(X)]^2=E(X^2)-[E(X)]^2.
\]
\end{infoblock}
\end{frame}

\begin{frame}[t]{Soru 17: Çözüm İçin Durak}
\small
\solutionghost{
\textbf{Kurulum}
\[
f(x)=
\begin{cases}
\dfrac{x}{2}, & 0<x<2,\\
0, & \text{diğer}.
\end{cases}
\]

\textbf{Kendinize sorun}
\begin{itemize}
  \item \(E(X)\) ve \(E(X^2)\) için hangi iki integral kurulmalı?
  \item Varyans farkı hangi iki sayının üzerinden bulunacak?
  \item \(g(X)=4X+3\) için lineerlik nasıl kullanılacak?
\end{itemize}
}
\end{frame}

\begin{frame}[t,shrink=2]{Soru 17: Çözüm}
\scriptsize
\solutionsetup{sürekli beklenen değer ve varyans}{moment integralleri ve lineerlik}{sürekli değişkende tüm ana nicelikler yoğunluğun integralleriyle elde edilir.}
\solutionmistake{\(E(X^2)\) ile \([E(X)]^2\) aynı nicelik değildir; varyans bu ikisinin farkından çıkar.}
\solutionstep{1}{Önce birinci ve ikinci momenti integralden hesaplıyoruz; çünkü varyans bunların farkına dayanır.}
\solutionmath{
E(X)=\frac12\int_0^2 x^2\,dx=\frac43,\qquad
E(X^2)=\frac12\int_0^2 x^3\,dx=2.
}

\solutionstep{2}{Varyans ve standart sapma doğrudan buradan gelir.}
\solutionmath{
\mathrm{Var}(X)=2-\left(\frac43\right)^2=\frac29,\qquad
\sigma_X=\frac{\sqrt2}{3}.
}

\solutionstep{3}{Lineerlik ile dönüşmüş beklenen değeri buluyoruz.}
\solutionmath{
E(4X+3)=4E(X)+3=\frac{25}{3}.
}
\solutionfinal{Bu yoğunlukta ortalama destek aralığının sağ tarafına kaymıştır; yani daha büyük \(x\) değerleri küçük değerlere göre biraz daha baskındır.}
\end{frame}

\begin{frame}[t]{Soru 18: Elektronik Bileşen Ömrü}
\footnotesize
Bir elektronik bileşenin ömrü \(X\) saat cinsinden ölçülüyor ve yoğunluk fonksiyonu aşağıdaki gibi veriliyor:
\[
f(x)=
\begin{cases}
\dfrac{20000}{x^3}, & x>100,\\
0, & \text{diğer}.
\end{cases}
\]

\textbf{İstenen}
\begin{itemize}
  \item[A.] Geçerli PDF olduğunu gösteriniz.
  \item[B.] Beklenen ömrü bulunuz.
  \item[C.] \(P(X>200)\) olasılığını bulunuz.
  \item[D.] \(E(1/X)\) değerini bulunuz.
\end{itemize}

\begin{infoblock}{Hatırlatma}
Sürekli durumda
\[
E(X)=\int_{-\infty}^{\infty} x\,f(x)\,dx,\qquad
E\!\left(g(X)\right)=\int_{-\infty}^{\infty} g(x)\,f(x)\,dx.
\]
\end{infoblock}
\end{frame}

\begin{frame}[t]{Soru 18: Çözüm İçin Durak}
\small
\solutionghost{
\textbf{Kurulum}
\[
f(x)=
\begin{cases}
\dfrac{20000}{x^3}, & x>100,\\
0, & \text{diğer}.
\end{cases}
\]

\textbf{Kendinize sorun}
\begin{itemize}
  \item Önce toplam alanın \(1\) olduğunu hangi uygunsuz integral ile göstereceksiniz?
  \item Beklenen ömür için integrand neden \(x\cdot f(x)\) olur?
  \item \(E(1/X)\) hesabında hangi yeni integrand ortaya çıkar?
\end{itemize}
}
\end{frame}

\begin{frame}[t,shrink=6]{Soru 18: Çözüm}
\scriptsize
\solutionsetup{uygunsuz integral ile PDF ve beklenen değer}{normlama, kuyruk olasılığı, dönüşmüş beklenen değer}{destek sonsuza gittiği için tüm hesaplar uygunsuz integral mantığıyla yapılır.}
\solutionmistake{Sonsuz üst sınırda uygunsuz integral mantığını unutmayın; hesap limit yorumuyla tamamlanır.}
\solutionstep{1}{Önce toplam alanı kontrol ediyoruz; çünkü sonsuz destekli bir yoğunlukta geçerlilik ayrıca doğrulanmalıdır.}
\solutionmath{
\int_{100}^{\infty}\frac{20000}{x^3}\,dx
=20000\left[-\frac{1}{2x^2}\right]_{100}^{\infty}
=1.
}

\solutionstep{2}{Beklenen ömrü ve kuyruk olasılığını hesaplıyoruz.}
\solutionmath{
E(X)=\int_{100}^{\infty}x\frac{20000}{x^3}\,dx
=20000\int_{100}^{\infty}x^{-2}\,dx
=200.
}
\solutionmath{
P(X>200)=\int_{200}^{\infty}\frac{20000}{x^3}\,dx=\frac14.
}

\solutionstep{3}{Dönüşmüş beklenen değeri aynı yöntemle buluyoruz.}
\solutionmath{
E\!\left(\frac1X\right)=\int_{100}^{\infty}\frac1x\frac{20000}{x^3}\,dx
=20000\int_{100}^{\infty}x^{-4}\,dx
=\frac{1}{150}.
}
\solutionfinal{Bu modelde bileşenin ortalama ömrü \(200\) saattir; yani ağır kuyruklu görünmesine rağmen sistemin tipik çalışma süresi hâlâ sonlu ve yorumlanabilirdir.}
\end{frame}

\begin{frame}[t]{Soru 19: \(a+bx\) Yoğunluğu}
\footnotesize
Sürekli bir rassal değişken \(X\) için yoğunluk fonksiyonu lineer biçimde veriliyor:
\[
f(x)=
\begin{cases}
a+bx, & 0<x<2,\\
0, & \text{diğer}
\end{cases}
\qquad\text{ve}\qquad
E(X)=\frac76.
\]

\textbf{İstenen}
\begin{itemize}
  \item[A.] \(a\) ve \(b\) bulunuz.
  \item[B.] Negatif olmama koşulunu doğrulayınız.
  \item[C.] \(P(X>1)\) bulunuz.
  \item[D.] \(\mathrm{Var}(X)\) hesaplayınız.
\end{itemize}

\begin{infoblock}{Hatırlatma}
Sürekli durumda
\[
E(X)=\int_{-\infty}^{\infty} x\,f(x)\,dx,\qquad
E(X^2)=\int_{-\infty}^{\infty} x^2 f(x)\,dx.
\]
\[
\mathrm{Var}(X)=E(X^2)-[E(X)]^2.
\]
\end{infoblock}
\end{frame}

\begin{frame}[t]{Soru 19: Çözüm İçin Durak}
\small
\solutionghost{
\textbf{Kurulum}
\[
f(x)=a+bx,\qquad 0<x<2,\qquad E(X)=\frac76.
\]

\textbf{Kendinize sorun}
\begin{itemize}
  \item Normlama integrali hangi ilk denklemi verir?
  \item Verilen ortalama bilgisi hangi ikinci denklemi üretir?
  \item \(a\) ve \(b\) bulunduktan sonra \(P(X>1)\) ve varyans nasıl tamamlanır?
\end{itemize}
}
\end{frame}

\begin{frame}[t,shrink=2]{Soru 19: Çözüm}
\scriptsize
\solutionsetup{parametreli sürekli yoğunluk ve moment koşulu}{normlama integrali, ortalama denklemi, ikinci moment}{iki bilinmeyen parametreyi çözmek için iki bağımsız koşul gerekir.}
\solutionmistake{Normlama denklemi ile ortalama denklemini karıştırmayın; biri toplam alanı, diğeri merkez bilgisini verir.}
\solutionstep{1}{Normlama ve verilen ortalama bilgisi bize \(a\) ve \(b\) için iki denklem verir.}
\solutionmath{
a+b=\frac12,\qquad
2a+\frac{8b}{3}=\frac76.
}

\solutionstep{2}{Doğrusal sistemi çözüyoruz.}
\solutionmath{
a=\frac14,\qquad b=\frac14,\qquad f(x)=\frac{1+x}{4}.
}

\solutionstep{3}{Olasılık ve varyansı tamamlıyoruz.}
\solutionmath{
P(X>1)=\int_1^2 \frac{1+x}{4}\,dx=\frac58,\qquad
E(X^2)=\frac53,\qquad
\mathrm{Var}(X)=\frac{11}{36}.
}
\solutionfinal{Bu modelde yoğunluk sağa doğru arttığı için \(X\)'in büyük değerlere kayma eğilimi vardır; bu da \(P(X>1)\) olasılığını yarının üstüne taşır.}
\end{frame}

\begin{frame}[t]{Ek Zor Soru: Firma Karşılaştırması}
\footnotesize
Firma A:
\[
\begin{array}{c|ccc}
x & 1 & 2 & 3\\ \hline
f_A(x) & 0.3 & 0.4 & 0.3
\end{array}
\]

Firma B:
\[
\begin{array}{c|ccccc}
x & 0 & 1 & 2 & 3 & 4\\ \hline
f_B(x) & 0.2 & 0.1 & 0.3 & 0.3 & 0.1
\end{array}
\]

\textbf{İstenen}: ortalamaları, varyansları ve daha kararsız firmayı bulunuz.
\end{frame}

\begin{frame}[t]{Ek Zor Soru: Çözüm İçin Durak}
\small
\solutionghost{
\textbf{Kurulum}
\[
\text{İki firma için de önce }E(X)\text{ ve }E(X^2)\text{ ayrı ayrı hesaplanacak.}
\]

\textbf{Kendinize sorun}
\begin{itemize}
  \item Ortalamalar eşit çıkarsa karar vermek için hangi ölçüye bakacaksınız?
  \item “Daha kararsız firma” ifadesi varyans açısından ne demektir?
\end{itemize}
}
\end{frame}

\begin{frame}[t,shrink=3]{Ek Zor Soru: Çözüm}
\footnotesize
\solutionsetup{iki dağılımı karşılaştırma}{beklenen değer, ikinci moment, varyans}{aynı ortalama her zaman aynı risk veya yayılım anlamına gelmez.}
\solutionmistake{İki dağılımı yalnızca ortalamayla karşılaştırmayın; kararlılık için varyans da gerekir.}
\solutionstep{1}{Önce iki firmanın ortalamalarını hesaplıyoruz; çünkü merkez aynı mı farklı mı önce bunu görmek gerekir.}
\solutionmath{
E_A(X)=1(0.3)+2(0.4)+3(0.3)=2,\qquad
E_B(X)=0(0.2)+1(0.1)+2(0.3)+3(0.3)+4(0.1)=2.
}

\solutionstep{2}{İkinci momentlerden varyansları çıkarıyoruz.}
\solutionmath{
E_A(X^2)=1(0.3)+4(0.4)+9(0.3)=4.6,\qquad
\mathrm{Var}_A(X)=4.6-2^2=0.6.
}
\solutionmath{
E_B(X^2)=0+1(0.1)+4(0.3)+9(0.3)+16(0.1)=5.6,\qquad
\mathrm{Var}_B(X)=5.6-2^2=1.6.
}
\solutionfinal{İki firmanın ortalama performansı aynı olsa da Firma B’nin sonuçları daha geniş saçıldığı için operasyonel açıdan daha kararsız olan firma B’dir.}
\end{frame}

\section{Kapanış}
\subsection{Özet}

\begin{frame}[t]{Bu Zor Setten Ne Çıkmalı?}
\small
\begin{itemize}
  \item İlk karar: soru olay mı soruyor, dağılım mı soruyor? Olaydaysa küme diliyle, dağılımdaysa PMF/PDF/CDF diliyle başlayın.
  \item Koşullu olasılıkta önce ``hangi bilgi örnek uzayı daralttı?'' sorusunu sorun; Bayes'te ise gözlenen olayı normalize edip yönü ters çevirin.
  \item CDF ayrık ise sıçrama büyüklüğü tek nokta olasılığı verir; sürekli ise tek nokta değil aralık alanı anlam taşır.
  \item Beklenen değer merkez bilgisini, varyans ise yayılımı verir; bu yüzden karşılaştırma sorularında yalnız ortalamaya güvenmeyin.
\end{itemize}

\begin{resultblock}{Son mesaj}
Bu setin ana kazanımı tek tek formül ezberlemek değil, doğru soruya doğru aracı seçebilmektir.
\end{resultblock}
\end{frame}

\begin{frame}[t]{Çıkış Bileti: Hangi Soruda Hangi Yöntem?}
\small
\begin{itemize}
  \item Bir makineden mi geldi, bir hastalık var mı, bir test pozitif mi? Böyle bir ters yön sorusunda neden önce toplam olasılık, sonra Bayes gerekir?
  \item Size \(F(x)\) verilirse hangi durumda ``sıçrama farkı'' ile \(P(X=x)\) bulursunuz, hangi durumda tek nokta olasılığının sıfır olduğunu söylersiniz?
  \item İki dağılımın ortalaması eşitse ``aynı davranıyorlar'' demek neden hatalıdır; hangi ikinci nicelik kararı değiştirir?
  \item Bir çözümün ilk satırında artık şu ayrımı net yapabiliyor musunuz: PMF mi, PDF mi, CDF mi, yoksa moment hesabı mı?
\end{itemize}
\end{frame}

\section{Kaynaklar}
\begin{frame}[allowframebreaks]{Kaynaklar}
  \nocite{walpole}
  \nocite{walpole_ch2}
  \nocite{walpole_ch3}
  \nocite{walpole_ch4}
  \printbibliography{}
\end{frame}

\appendix

\section{Yapay Zeka Asistanı}
\begin{frame}{Yapay Zeka Asistanınız: NotebookLM}
  \begin{columns}[T]
    \begin{column}{0.48\textwidth}
      \justifying
      Bu deck'teki zor soruları küçük parçalara ayırıp tekrar etmek için
      NotebookLM kullanılabilir.

      \vspace{0.8em}
      \begin{itemize}
        \item soru bazlı tekrar,
        \item kendi çözümünüzle karşılaştırma,
        \item kısa özet üretme,
        \item bölüm bazlı çalışma listesi çıkarma.
      \end{itemize}
    \end{column}
    \begin{column}{0.48\textwidth}
      \centering
      \includegraphics[width=0.82\linewidth]{figures/notebooklm.jpg}
    \end{column}
  \end{columns}
\end{frame}

\begin{frame}[plain,noframenumbering]
  \begin{tikzpicture}[remember picture,overlay]
    \node[anchor=center, at=(current page.center), inner sep=0pt] {
      \includegraphics[width=\paperwidth, height=\paperheight]{logos/background_figure.jpg}
    };
  \end{tikzpicture}
\end{frame}

\end{document}
